\section{Opening the Liturgical Year}
\runninghead{Opening the Liturgical Year}

On the First Sunday of Advent, the Roman Church opens her new liturgical year at the Basilica of St.~Mary Major on the Esquiline Hill (\latin{Statio ad S.~Mariam maiorem}, ``Station at St.~Mary Major''), the Roman church of the Crib, wearing violet vestments and keeping silent the angels' Christmas hymn \latin{Gloria in excelsis Deo} (``Glory to God in the highest''). The Gospel appointed for this opening Mass is not the Annunciation or the journey to Bethlehem, but Christ's warning from the Temple courts during Passion Week: signs in the sun, moon, and stars, nations in distress at the roaring of the sea, men withering in terror, and the Son of Man coming in a cloud with great power and majesty (Luke~21:25--33). Yet in the very same liturgy, the Epistle calls the sleeping assembly to put on the armour of light because day is at hand (Rom.~13:11--14a), the Alleluia begs God to show His mercy and give His salvation (Ps.~84:8), and the Communion antiphon announces that the Lord will give goodness and our earth shall yield her fruit (Ps.~84:13).

How do cosmic collapse, a call to wake at dawn, a Marian stational basilica, and a promise of fruitful soil fit into one Sunday's worship? The Church Fathers, Doctors, and liturgical commentators answer that question along three distinct lines of interpretation. Along the first, St.~Gregory the Great expounds Luke~21:25--33 as a summons to detach from an aging world before the strict Judge, St.~Cyril of Alexandria explains its promised redemption as bodily resurrection and cosmic restoration, St.~Bede (on Luke~21:33) and St.~Augustine (\work{De civitate Dei}~XX) clarify how heaven and earth pass away in outward form while remaining in nature, St.~John Chrysostom expounds Romans~13:11--14 as a dawn call to battle-array before the approaching Judgment, and St.~Robert Bellarmine expounds Psalm~24:1--3 as tearing the soul free from earthly briers. Along the second, St.~Ambrose of Milan interprets the cloud of Luke~21:27 as the prophets and the Virgin Mary overshadowed by the Holy Spirit and the sprouting trees of verse~29 as Israel and all the nations blossoming toward resurrection, St.~Robert Bellarmine hears in Psalm~84:8 and~13 the Incarnation of the Son through the Holy Spirit overshadowing Mary's virgin earth, St.~Augustine (\work{Enarrationes in Psalmos}~84.13 and \work{Sermo}~185.1) identifies Truth sprung from earth in Psalm~84:12 with the Son born of Virgin flesh, and St.~Bernard of Clairvaux preaches Christ as the Fruit of Nazareth (\work{Missus est}, Hom.~1) and Mary's womb as an unploughed field (\work{Sermones in Adventu}~2). Along the third, St.~Augustine expounds Psalm~84:8 and~13 (working from his Old Latin text reading \latin{suavitatem}, ``sweetness'') as prevenient grace furrowing the sinner's heart with confession, St.~Thomas Aquinas traces Psalm~24:1--4 from contemplative ascent through the highway of commandments to the narrower shortcut of the evangelical counsels, and Bl.~Ildefonso Schuster expounds the Secret, Communion, and Postcommunion as the Eucharistic Mysteries purifying worshippers, refreshing arid hearts, and preparing the Church for Christmas.

Under the universal General Roman Calendar of the 1962 \work{Missale Romanum}, \latin{Dominica prima Adventus} (``First Sunday of Advent'') is a Sunday of the first class (\latin{I~classis}, \latin{Rubricae generales} [RG]~11, 15, and Table of Precedence RG~91, row~6), falling on the Sunday on or nearest to 30~November (RG~20): in 2026, Sunday, 29~November, opening the 2026--2027 liturgical year. Because a first-class Sunday admits only privileged commemorations (RG~109, 111~a; \latin{Rubricae generales Missalis romani} [RGMR]~434), the ordinary commemoration of St.~Saturninus, Martyr, occurring on 29~November is omitted in Mass, and the pre-1955 seasonal second and third orations were abolished by the 1960 \work{Variationes} (II.13, IV.18). Mass therefore has one Collect, one Secret, and one Postcommunion (RGMR~434, 454, 457~d, 480--481, 505). The prayers at the foot of the altar and the \latin{Gloria Patri} (``Glory be to the Father'') at the Introit remain (RGMR~424--429); the \latin{Gloria in excelsis} is omitted throughout Advent \latin{de Tempore} (``of the temporal cycle,'' printed rubric, p.~1; RGMR~431--432); the \latin{Credo} (``Creed'') is said on Sundays (printed rubric, p.~2; RGMR~475--476); the Preface of the Most Holy Trinity is appointed for Sundays of Advent (printed rubric, p.~2; RGMR~494~b), whereas Advent weekdays resuming the Sunday Mass omit the Alleluia and \latin{Credo} and take the Common Preface (printed rubrics, p.~2; RGMR~299, 482, 498). Under RGMR~507, Mass concludes with \latin{Ite, missa est} (``Go, the Mass is ended'') and the ordinary Last Gospel, John~1:1--14 (RGMR~509--510); dalmatic and tunicle are worn by sacred ministers, folded chasubles having been abolished in 1960 (RG~137).

For the composition dates, traditional attributions, historical settings, and narrated events of the seven appointed scriptural passages, see the appendix \hyperref[sec:date-location]{\work{Scriptural Date and Location}}. Psalm citations throughout follow the Missal's Vulgate numbering first, with Hebrew numbering in parentheses where it differs.

\subsection{Map of the Appointed Propers}

\begin{maptable}
\textbf{Introit} & \latin{Ad te levavi} (``To thee have I lifted up my soul''); Ps.~24:1--3, 4b (Heb.~25:1--3a, 4) & Lifts the soul to God in unashamed trust, confesses that none who wait on Him (\latin{qui te exspectant}) will be confounded, and asks to be shown His ways and taught His paths.\\
\textbf{Collect} & \latin{Excita, quaesumus, Domine} (``Bestir, O Lord, Thy might, we pray Thee''); composed prayer echoing Ps.~79:3 & Addressed directly to Christ (\latin{Qui vivis}, ``Who livest''): begs Him to rouse His power and come so that, protected and freed by Him, we may be rescued from the pressing dangers of our sins.\\
\textbf{Epistle} & \latin{Fratres: Scientes quia hora est} (``Brethren: Knowing that it is now the hour''); Rom.~13:11--14a & Summons the assembly to rise from sleep because salvation is nearer; bids them cast off works of darkness, put on the armour of light, walk decently as in daylight, and put on the Lord Jesus Christ.\\
\textbf{Gradual} & \latin{Universi, qui te exspectant} (``All who wait on thee''); Ps.~24:3b, 4b (Heb.~25:3a, 4) & Repeats the promise that none who wait on the Lord will be confounded and asks Him in the Old Latin wording (\latin{notas fac mihi}, ``make known to me'') to make His ways known and teach His paths.\\
\textbf{Alleluia} & \latin{Ostende nobis, Domine} (``Shew us, O Lord, thy mercy''); Ps.~84:8 (Heb.~85:8; omitted on resumed Advent weekdays) & Sings \latin{Alleluia} on Sunday even while \latin{Gloria} is silent, pleading that the Lord show His mercy and grant His salvation (\latin{salutare tuum}).\\
\textbf{Gospel} & \latin{Erunt signa in sole} (``There shall be signs in the sun''); Luke~21:25--33 & Foretells celestial signs, the distress of nations, and the Son of Man coming in a cloud with power and majesty; commands disciples to lift up their heads, learn from the fig tree and all the trees, and trust Christ's unfailing words.\\
\textbf{Offertory} & \latin{Ad te levavi} (``To thee have I lifted up my soul''); Ps.~24:1--3 (Heb.~25:1--3a) & Repeats the Introit antiphon verbatim at the preparation of the gifts, entrusting the worshipping community to God before the sacrifice.\\
\textbf{Secret} & \latin{Haec sacra nos, Domine} (``Cleansing us by their mighty power''); composed prayer & Prays that these sacred Mysteries, having cleansed us by their mighty power, may bring us purer to their own Source (\latin{principium}).\\
\textbf{Communion} & \latin{Dominus dabit benignitatem} (``The Lord will give goodness''); Ps.~84:13 (Heb.~85:13) & Proclaims that the Lord will give goodness (\latin{benignitatem}) and our earth shall yield her fruit.\\
\textbf{Postcommunion} & \latin{Suscipiamus, Domine, misericordiam tuam} (``May we receive Thy mercy, O Lord''); composed prayer echoing Ps.~47:10 & Prays that we may receive God's mercy in the midst of His temple so as to go before the coming solemnities of our reparation with fitting honour.\\
\end{maptable}
